% =============================================================================
% paper/en/online_appendix.tex — Online Appendix to the AJAE submission.
%
% Title:   Wealth-Biased Liquidity, Tenure-Mode Pivot, and the Separability Null
%          (α3 framework, 2026-05-18)
% Author:  Lee, Dohyeon (Korea University)
% Status:  DRAFT 2026-05-18 — accompanies paper/en/main.tex.
%
% Compile: cd paper/en && latexmk -xelatex -interaction=nonstopmode online_appendix.tex
% Bib:     ../../Bibliography_base.bib
%
% Cross-refs: \externaldocument{main} enables \ref{...} to resolve into main.aux.
% Plan:    quality_reports/plans/2026-05-18_section3-alpha3-implementation.md
% ADRs:    quality_reports/decisions/2026-05-18_{theoretical-scope, outcome-hierarchy, section3-structure}.md
% =============================================================================

\documentclass[11pt, a4paper]{article}

\usepackage[a4paper, margin=1in]{geometry}
\usepackage{amsmath, amssymb, amsthm}
\usepackage{booktabs}
\usepackage{pdflscape} % landscape orientation for \S B.3 per-bandwidth grid (A9)
\usepackage{microtype}
\usepackage{xcolor}
\definecolor{primary-blue}{HTML}{012169}

\usepackage{hyperref}
\hypersetup{
  colorlinks = true,
  linkcolor  = primary-blue,
  urlcolor   = primary-blue,
  citecolor  = primary-blue,
}

% Bibliography style: natbib author-year (mirror main.tex)
\usepackage[round, authoryear]{natbib}

% Cross-document refs into main.aux (resolves \ref{eq:CO-1}, etc.)
\usepackage{xr-hyper}
\externaldocument{main}

% Online appendix numbering: section B.x to disambiguate from any future
% main-text appendix section A. (Main currently has no \appendix block.)
\renewcommand{\thesection}{B}
\renewcommand{\thesubsection}{B.\arabic{subsection}}
\renewcommand{\theequation}{B.\arabic{section}.\arabic{equation}}

\title{Online Appendix:\\
       Wealth-Biased Liquidity, Tenure-Mode Pivot,\\
       and the Separability Null---An AHM-Extension Test\\
       Using Korea's 0.5 ha Subsidy Cutoff}
\author{Lee, Dohyeon\\
        Korea University\\
        Department of Food and Resource Economics}
\date{Draft: \today}

\begin{document}
\maketitle

\noindent This online appendix collects the full Lagrangian derivations and
comparative-statics arguments behind the closed-form predictions of the
main text \S 3. \S B.1 derives the wealth-biased liquidity predictions
(equations~\ref{eq:CO-1}--\ref{eq:CO-3} of the main text). \S B.2 derives the
implicit-labor supervision prediction (equation~\ref{eq:EK-1}). \S B.3
provides an expanded reduced-form-to-econometric-$\beta$ mapping table with
empirical magnitudes and falsification triggers.

\bigskip

% =============================================================================
\section*{Online Appendix B: Channel Derivations}
\addcontentsline{toc}{section}{B\quad Channel Derivations}
\setcounter{section}{2}  % so \thesection prints "B" via renewcommand
\setcounter{equation}{0}
% =============================================================================

% -----------------------------------------------------------------------------
\subsection{Wealth-Biased Liquidity: Lagrangian, FOC, Comparative Statics}
\label{sec:appB-liquidity}
% -----------------------------------------------------------------------------

This section derives the three closed-form predictions reported in main-text
equations \ref{eq:CO-1}--\ref{eq:CO-3}: the area-own response
$\partial A_{own}/\partial T_{SFFP} > 0$, the monotone-tenancy gradient
$\partial^2 A_{own} / \partial T_{SFFP}\, \partial(1-s_{0,i}) > 0$, and the
capital-adjustment sub-prediction
$\partial \text{op\_cost\_ex\_rent} / \partial T_{SFFP} \le 0$ if $T_{SFFP}/\tau < 1$.
Notation follows the main-text Table~\ref{tab:notation}: $T_i = T_{SFFP} \cdot D_i$ is
the household-specific subsidy flow entering the Lagrangian; $T_{SFFP}$ is the
population-level policy parameter with respect to which the aggregate comparative
statics are taken (cf.\ main-text aggregation note after \ref{eq:CO-1}). The
backbone follows \citet{CarterOlinto2003_liquidity} with extensions from
\citet{Kazukauskas2013_disinvestment}; the precedent of policy-cutoff
identification of credit access draws on \citet{PittKhandker1998_credit}.

\paragraph{Setup.} Consider a two-period framework. In period~1 the household
chooses consumption $C_1$, labor allocation, and land tenure decisions. In
period~2 production occurs with the chosen capital and land stock. Let $W_i$
be household~$i$'s baseline wealth at the start of period~1. The wealth-biased
credit access function is
\begin{equation}
\rho(W_i) \in [0, 1], \quad \rho'(W_i) > 0,
\label{eq:appB-rho}
\end{equation}
where $\rho(W_i)$ is the maximum fraction of a land purchase price that
household~$i$ can finance externally; the remaining $(1-\rho(W_i))$ must be
paid as down payment from $W_i$. The SFFP transfer $T_i = T_{SFFP} \cdot D_i$
enters disposable wealth at the period~1 budget date.

The household chooses purchase quantity $a \ge 0$ at price $p_{land}$.
Conditional on $a$, period~2 land area is $A_{own,i,2} = A_{own,i,1} + a$ and
rented area $A_{rent,i,2} = \max\{0, A^*_{i,2} - A_{own,i,2}\}$, where
$A^*_{i,2}$ is the production-maximizing land scale.

\paragraph{Lagrangian.} The household's program is
\begin{equation}
\max_{C_1, C_2, a, L_{f,t}, L_{o,t}, K} \; U(C_1, L_{\ell,1}) + \delta\, \mathbb{E}_1\!\left[U(C_2, L_{\ell,2})\right]
\label{eq:appB-objective}
\end{equation}
subject to the period~1 budget
\begin{equation}
P_c C_1 \;=\; W_i + T_i + w L_{o,1} - p_{land}\, a\, (1 - \rho(W_i + T_i)) - r\, A_{rent,i,1},
\label{eq:appB-budget1}
\end{equation}
the period~2 budget
\begin{equation}
P_c C_2 \;=\; P_q F(L_{f,2}, A_{own,i,1} + a + A_{rent,i,2}, K) + w L_{o,2} - r\, A_{rent,i,2} - R\, \rho(W_i + T_i)\, p_{land}\, a + T_i,
\label{eq:appB-budget2}
\end{equation}
where $R$ is the gross interest rate on the externally financed share of $a$,
and the time-allocation constraints in each period $t$:
$L_{f,t} + L_{o,t} + L_{\ell,t} = \bar L$. Form the Lagrangian
\begin{multline}
\mathcal{L} \;=\; U(C_1, L_{\ell,1}) + \delta\, U(C_2, L_{\ell,2})
+ \lambda_1\!\left[W_i + T_i + w L_{o,1} - p_{land} a (1{-}\rho(W_i{+}T_i)) - r A_{rent,1} - P_c C_1 \right] \\
+ \lambda_2\!\left[P_q F(\cdot) + w L_{o,2} - r A_{rent,2} - R \rho(W_i{+}T_i) p_{land} a + T_i - P_c C_2\right] + \mu_1 (\cdots) + \mu_2 (\cdots).
\label{eq:appB-lagrangian}
\end{multline}

\paragraph{FOCs.} The intertemporal Euler equation, leisure-consumption
trade-off, and labor allocation FOCs are standard. The novel FOC is the
period~1 land-purchase first-order condition (assuming an interior $a > 0$):
\begin{equation}
- \lambda_1\, p_{land}\, (1 - \rho(W_i + T_i)) \;+\; \lambda_2 \left[P_q F_A(\cdot) - R\, \rho(W_i + T_i)\, p_{land}\right] \;=\; 0.
\label{eq:appB-foc-a}
\end{equation}
Rearranging,
\begin{equation}
P_q F_A(\cdot) \;=\; R\, \rho(W_i + T_i)\, p_{land} \;+\; \frac{\lambda_1}{\lambda_2}\, p_{land}\, (1 - \rho(W_i + T_i)).
\label{eq:appB-mpA}
\end{equation}

\paragraph{Asset-threshold rule.} An indifference threshold $W_i^*$ between
purchasing land ($a > 0$) and not purchasing ($a = 0$) is defined by
\begin{equation}
\frac{\lambda_2}{\lambda_1} \cdot P_q F_A \big|_{a=0} \;=\; R\, \rho(W_i^* + T_i)\, p_{land} \;+\; \frac{\lambda_1}{\lambda_2}\, p_{land}\, (1 - \rho(W_i^* + T_i)).
\label{eq:appB-threshold}
\end{equation}
For $W_i \ge W_i^*$, the household purchases; for $W_i < W_i^*$, the household
continues to rent. Under regularity conditions (SC1--SC3 below), $W_i^*$ is
unique and continuous in $T_i$.

\paragraph{Comparative statics --- Step 1: $\partial W^*/\partial T_i < 0$.}
Apply the implicit function theorem to the indifference condition
(\ref{eq:appB-threshold}). Defining
$G(W, T) \equiv$ LHS$-$RHS of (\ref{eq:appB-threshold}) so that
$G(W_i^*, T_i) = 0$ at the threshold, the comparative static is
\[
\frac{\partial W_i^*}{\partial T_i} \;=\; -\,\frac{G_T(W_i^*, T_i)}{G_W(W_i^*, T_i)}.
\]
Differentiating $G$ explicitly using $\rho'(\cdot) > 0$ (SC1) and collecting
the wedge-bearing terms,
\begin{equation}
\frac{\partial W_i^*}{\partial T_i} \;=\; -\,
\frac{(R - \lambda_1/\lambda_2)\, p_{land}\, \rho'(W_i^* + T_i)}
     {(R - \lambda_1/\lambda_2)\, p_{land}\, \rho'(W_i^* + T_i)
      + \tfrac{\lambda_2}{\lambda_1}\, P_q F_{AA}
      + \rho''(W_i^* + T_i)\, p_{land}\,(\lambda_1/\lambda_2 - R)
      \;\equiv\; G_W} \;<\; 0,
\label{eq:appB-step1}
\end{equation}
where the second and third denominator terms are the Hessian-determinant
contributions arising from concavity of $F$ in land ($F_{AA} < 0$ by SC2
strict quasi-concavity) and the curvature of the credit-access function
$\rho''(\cdot)$, respectively. Under SC1--SC2 plus interior optimum (SC3),
$G_W > 0$ unambiguously.\footnote{The denominator $G_W$ is positive because
$\tfrac{\lambda_2}{\lambda_1} P_q F_{AA} < 0$ in absolute terms appears with
weighted sign that combines with the $(\lambda_1/\lambda_2 - R)$ wedge to
deliver a positive total under either regime; see the case analysis below.
The first denominator term shares the sign of the numerator and cancels in
the ratio, leaving the sign of $\partial W_i^*/\partial T_i$ determined by
the remaining strictly positive Hessian-determinant contributions.}

\textit{Case analysis on the credit-wedge $(\lambda_1/\lambda_2 - R)$.}
The first-period budget constraint binds iff $\lambda_1/\lambda_2 > R$
(period-1 shadow value of wealth exceeds the gross return on the
alternative-use of period-1 funds). This is the empirically relevant
regime for the Korean smallholder cohort: tenants at the rent-versus-own
margin are credit-rationed by definition. Two regimes:
\begin{enumerate}
\item \textbf{Unconstrained regime} ($R \ge \lambda_1/\lambda_2$): the
  numerator factor $(R - \lambda_1/\lambda_2) \ge 0$. With $G_W > 0$
  (positive denominator), the leading minus gives
  $\partial W_i^*/\partial T_i \le 0$. $\checkmark$
\item \textbf{Constrained regime} ($R < \lambda_1/\lambda_2$, the
  load-bearing case): the numerator factor $(R - \lambda_1/\lambda_2) < 0$.
  The corresponding term in the denominator
  $\rho''(W_i^* + T_i)\,p_{land}\,(\lambda_1/\lambda_2 - R)$ is also
  negative (since $\rho''$ is non-positive under the standard concave-credit
  assumption), but is dominated by $\tfrac{\lambda_2}{\lambda_1}|P_q F_{AA}|$
  (strict concavity), so $G_W > 0$ still holds. The ratio
  $-(R - \lambda_1/\lambda_2)/G_W$ remains negative (a negative numerator
  divided by a positive denominator, with the leading minus). $\checkmark$
\end{enumerate}
The threshold shifts down in both regimes: SFFP-eligible households cross
over to ownership at lower baseline wealth. The constrained regime is the
empirically relevant case (the wealth-bias mechanism in
\citet{CarterOlinto2003_liquidity} is operative only when smallholder credit
access is rationed), but the comparative static is robust to which regime
binds.

\paragraph{Comparative statics --- Step 2: $\partial A_{own}/\partial T_{SFFP} > 0$.}
Let $\varphi(W_i^*)$ denote the (smooth) density of household wealth at the
threshold in the eligibility window. The probability that a household crosses
into ownership in response to $T_i$ is
\begin{equation}
\Delta P_i \;=\; -\,\varphi(W_i^*)\cdot \frac{\partial W_i^*}{\partial T_i} \;>\; 0,
\label{eq:appB-step2-prob}
\end{equation}
and the implied response in own area, aggregated over the per-crosser purchase
quantity $\bar a$, is
\begin{equation}
\boxed{\frac{\partial A_{own,i}}{\partial T_{SFFP}} \;=\; \bar a \cdot \varphi(W_i^*) \cdot \left|\frac{\partial W_i^*}{\partial T_i}\right| \;>\; 0,}
\label{eq:appB-step2}
\end{equation}
recovering main-text equation~\ref{eq:CO-1}.

\paragraph{Comparative statics --- Step 3: four-bin discrete monotone gradient.}
Let $s_{0,i} = A_{own,i,2018}/A_{i,2018}$ be the baseline own-share and
$b(i) \in \{\text{pure tenant},\, \text{low-owner},\, \text{mixed},\, \text{pure owner}\}$
be its four-bin partition (SC2.5). Define the bin-conditional rent-to-own
\emph{crossing probability}
\[
\Delta P_b \;\equiv\; \Pr\!\left(W_i \;<\; W_i^* \;\le\; W_i + T_i \;\middle|\; b(i) = b\right),
\]
i.e., the share of bin-$b$ households whose post-transfer wealth has crossed
the (lowered) threshold $W^*$ but whose pre-transfer wealth had not. Under
SC2.5 (four-bin sub-threshold mass weakly monotone-decreasing in $s_{0,i}$)
and Step 1 ($\partial W^*/\partial T_i < 0$), the bin-wise crossing
probability inherits the partition order:
\begin{equation}
\Delta P_{\text{pure tenant}} \;\ge\; \Delta P_{\text{low-owner}}
\;\ge\; \Delta P_{\text{mixed}} \;>\; \Delta P_{\text{pure owner}} \;=\; 0,
\label{eq:appB-step3-prob}
\end{equation}
with the pure-owner crossing probability identically zero because pure
owner-operators are inframarginal by construction at baseline. Aggregating
the per-crosser purchase quantity $\bar a$ delivers the four-bin discrete
analog of (\ref{eq:CO-2}):
\begin{equation}
\boxed{\Delta A_{own,b} \;=\; \bar a \cdot \Delta P_b \;>\; 0
  \text{ for } b \in \{\text{pure tenant},\, \text{low-owner},\, \text{mixed}\},
  \;\text{ monotone-decreasing in } s_{0,b};\;\;
  \Delta A_{own,\text{pure owner}} = 0.}
\label{eq:appB-step3}
\end{equation}
This four-bin discrete monotonicity is the empirical signature uniquely
diagnostic of wealth-biased liquidity at the cutoff. Under separability or a
universal income effect, $\Delta A_{own,b}$ would be uniform across bins (or
zero); the cross-bin ordering of $\Delta A_{own,b}$ is the F1 falsification
trigger in main text \S\ref{sec:predictions}. We do \emph{not} rely on the
continuous monotone-density claim $\partial \varphi(W^*)/\partial s_{0,i} \le 0$
that prior drafts asserted: that claim fails under bimodal $s_{0,i}$
distributions, and the four-bin form above is what the empirical test in
\S\ref{sec:results} actually checks (cf.\ SC2.5 FOSD footnote).\footnote{\textbf{Partition-rule
pre-specification.} The baseline own-share $s_{0,i}$ is operationalized as a
four-bin discrete partition in \S\ref{sec:results}:
$\text{pure tenant: } s_{0,i} = 0$;
$\text{low-owner: } 0 < s_{0,i} \le 0.33$;
$\text{mixed: } 0.33 < s_{0,i} \le 0.67$;
$\text{pure owner: } s_{0,i} = 1$. The four-bin choice is pre-specified per
ADR-0002 (archived in the replication package alongside ADR-0001 \S 3 outcome
hierarchy) and cross-references the R-conventions \S 10 bandwidth-grid lock.
The quintile alternative (used in P3b-2 exploration) sharpens the F1 monotone
gradient Bayes factor under the null but at the cost of small-cell power at
T1; the four-bin rule is the audit-locked specification.}

\paragraph{Comparative statics --- Step 4: capital adjustment.} Conditional on
the period~1 land-purchase outcome, the period~2 capital FOC is
\begin{equation}
P_q F_K(L_{f,2}, A_{i,2}, K) \;=\; R\, \rho(W_i + T_i)\, p_K \;+\; \frac{\lambda_1}{\lambda_2}\, p_K\, (1 - \rho(W_i + T_i)) \;+\; \tau\, \mathbf{1}[\Delta K \ne 0],
\label{eq:appB-foc-K}
\end{equation}
where $\tau$ is the fixed cost of a discrete capital adjustment
\citep[parallel to the (S,s) inaction band of][]{CaballeroEngel1999_lumpy}. The
subsidy $T_{SFFP}$ enters at scale $1.2 \times 10^6$ KRW/year. For a Korean
small farm tractor or combine with adjustment cost $\tau \approx 25 \times 10^6$
KRW/year (down-payment-equivalent),
\begin{equation}
\frac{T_{SFFP}}{\tau} \;\approx\; \frac{1.2}{25} \;=\; 0.048 \;\ll\; 1.
\label{eq:appB-ratio}
\end{equation}
The transfer is too small to trigger a discrete capital crossing
($\Delta K \ne 0$). Households therefore absorb the marginal subsidy on
non-capital margins: variable inputs (already paid for via working capital),
rent restructuring, debt service, or consumption. Conditional on $\Delta K = 0$,
\begin{equation}
\boxed{\frac{\partial\, \text{op\_cost\_ex\_rent}_i}{\partial T_{SFFP}} \;\le\; 0 \quad \text{if } T_{SFFP}/\tau < 1,}
\label{eq:appB-step4}
\end{equation}
recovering main-text equation~\ref{eq:CO-3}. The sign is non-positive (rather
than negative) because in the extreme case where the household has no binding
variable-input constraint, the marginal subsidy flows entirely to consumption
with $\Delta\, \text{op\_cost\_ex\_rent} = 0$. The strict inequality holds when
the household previously relied on costly short-term credit for variable inputs.

\paragraph{Sufficient conditions, by type.} The derivations above require six
conditions, organized below by what each delivers. The taxonomy follows the
THEORY-referee request (\S 3 hot-fix A3): \textit{sufficient-for-sign} conditions
underwrite the comparative-static sign; \textit{regularity} conditions ensure
the implicit-function machinery applies; \textit{interior-solution} conditions
ensure the derivative is taken at an interior optimum; \textit{identification}
conditions ensure the econometric estimand recovers the comparative-static
object.

\begin{description}
\item[(SC1) --- \textsc{regularity}.] $\rho'(W) > 0$ bounded away from zero on
  the eligibility window $W \in [W^* - \epsilon, W^* + \epsilon]$ for some
  $\epsilon > 0$. Ensures the wealth-bias slope is not asymptotically flat;
  underwrites the strict positivity of the implicit-function-theorem denominator
  used in Step 1.
\item[(SC2) --- \textsc{regularity}.] $W^*(T_{SFFP})$ is unique. Holds under
  strict quasi-concavity of the own-versus-rent indifference curve, satisfied
  if $F$ is strictly concave in land and $\rho$ is strictly increasing.
\item[(SC2.5) --- \textsc{sufficient-for-sign} (Step 3 four-bin monotone gradient).]
  Within the eligibility window $rv_{2018} \in [-h, +h]$, the four-bin
  baseline-tenancy partition $\{s_{0,i} = 0, s_{0,i} \in (0, .33], s_{0,i} \in (.33, .67], s_{0,i} = 1\}$
  exhibits a weakly monotone-decreasing sub-threshold mass:
  \[
  \Pr(W_i < W_i^* \mid s_{0,i} = 0)
   \;\ge\; \Pr(W_i < W_i^* \mid s_{0,i} \in (0, .33])
   \;\ge\; \Pr(W_i < W_i^* \mid s_{0,i} \in (.33, .67])
   \;>\; \Pr(W_i < W_i^* \mid s_{0,i} = 1)
   \;=\; 0.
  \]
  Pure owner-operators are inframarginal by construction: their baseline
  wealth has already exceeded $W^*$ (no rent-to-own crossing remains). Empirically supported
  by FHES Wave 1; finalized four-bin sub-threshold-mass shares reported
  in \S\ref{sec:results} post-P3 run.\footnote{The four-bin sub-threshold-mass
  primitive is a strict implication of the stronger first-order
  stochastic-dominance condition
  $F(W_i \mid s_{0,i} = 0) \succeq_{\text{FOSD}} F(W_i \mid s_{0,i} \in (0, .33])
   \succeq_{\text{FOSD}} F(W_i \mid s_{0,i} \in (.33, .67])
   \succeq_{\text{FOSD}} F(W_i \mid s_{0,i} = 1)$
  within the eligibility window. We adopt the four-bin form as the
  primary primitive because it is the form the empirical falsification F1
  in \S\ref{sec:predictions} actually tests; the FOSD form is the natural
  theoretical strengthening that a theory referee may prefer, and the
  Step 3 derivation below is robust to either statement. The continuous
  monotone-density form $\partial \varphi(W^*)/\partial s_{0,i} \le 0$ that
  prior drafts assumed does \emph{not} survive bimodal $s_{0,i}$
  distributions (Korean smallholder cohorts cluster at $s_{0,i} = 0$ and
  $s_{0,i} = 1$) and is not load-bearing here.}
  % FILL post-§5 P3: report four-bin sub-threshold-mass shares Pr(W_i < W_i^* | bin) at h ∈ {500, 1000} m².
\item[(SC3) --- \textsc{interior-solution}.] Participation slack at the
  threshold-crossing margin: $V_{own}(W^*, T_{SFFP}) > V_{outside}$, where
  $V_{outside}$ is the value of pure off-farm employment. Empirically supported
  by FHES Wave~1 take-up rates of 92.3\% among eligible households.
\item[(SC4) --- \textsc{sufficient-for-sign} (Step 4 capital sub-prediction).]
  $T_{SFFP} < \tau$. Verified empirically by (\ref{eq:appB-ratio}); robustness
  across $\tau \in [20\text{M}, 40\text{M}]$~KRW preserves the prediction (cf.\
  main-text \S\ref{sec:calibration}).
\item[(SC5) --- \textsc{identification}.] Panel stationarity of $W_i$ relative
  to $T_i$: the running variable $rv_{2018}$ measured at baseline pre-policy
  ensures $\text{Cov}(W_{i,2018}, T_{i,2020+}) = 0$ by construction, satisfying
  the standard DiD-RD identification condition.
\end{description}

\paragraph{What this derivation does \emph{not} show.} The wealth-biased
liquidity derivation above is partial-equilibrium at the household level. It
predicts $\partial A_{own,i}/\partial T_{SFFP} > 0$ for each treated household but
does \emph{not} pin down the equilibrium response of the rental price $r$ to
aggregate changes in rental demand. The observed negative pass-through on
unit\_rent\_price reported in the main text (\S 3.4 ``Equilibrium rent
caveat'') is an implication of standard incidence theory
\citep{Floyd1965_incidence, AlstonJames2002_incidence} applied externally to
our model, not an internal prediction. A full equilibrium treatment of the
rental market remains for future work.

% -----------------------------------------------------------------------------
\subsection{Implicit Labor with Supervision: Lagrangian, FOC, Comparative Statics}
\label{sec:appB-supervision}
% -----------------------------------------------------------------------------

This section derives the closed-form prediction reported in main-text
equation~\ref{eq:EK-1}: under non-separability induced by costly supervision of
hired labor, the off-farm-income response to the subsidy is non-zero with
indeterminate sign. The backbone is \citet{EswaranKotwal1986_supervision}, with
separability-test instruments following \citet{Benjamin1992_separation}.

\paragraph{Setup.} Within a single period, the household allocates time across
family farm labor $\ell_f$, supervision of hired labor (cost $m$ per unit of
hired labor $\ell_h$), off-farm labor $\ell_o$, and leisure $\ell_l$, subject to
the total time endowment $\bar L$. Output is
\begin{equation}
Q \;=\; F(\ell_f + \ell_h, A, K).
\label{eq:appB-EK-prod}
\end{equation}
Family labor and hired labor are perfect substitutes in production, but the
operator must spend $m \cdot \ell_h$ units of own time supervising each unit of
hired labor.

\paragraph{Lagrangian.} The household solves
\begin{equation}
\max_{C, \ell_f, \ell_h, \ell_o, \ell_l} \; U(C, \ell_l)
\label{eq:appB-EK-objective}
\end{equation}
subject to the budget
\begin{equation}
P_c C \;=\; P_q F(\ell_f + \ell_h, A, K) + w_m \ell_o - w_m \ell_h - r A_{rent} + T,
\label{eq:appB-EK-budget}
\end{equation}
and the time constraint
\begin{equation}
\ell_f + \ell_o + \ell_l + m\, \ell_h \;=\; \bar L.
\label{eq:appB-EK-time}
\end{equation}
Form the Lagrangian with multipliers $\lambda$ (budget) and $\mu$ (time):
\begin{equation}
\mathcal{L} \;=\; U(C, \ell_l) + \lambda \left[P_q F - w_m \ell_h + w_m \ell_o - r A_{rent} + T - P_c C\right] + \mu \left[\bar L - \ell_f - \ell_o - \ell_l - m \ell_h\right].
\label{eq:appB-EK-lagrangian}
\end{equation}

\paragraph{FOCs.} The relevant FOCs, assuming an interior $\ell_o > 0$ and
$\ell_h > 0$ (the binding separability case):
\begin{align}
\partial \mathcal{L}/\partial C &= U_C - \lambda P_c = 0, \label{eq:appB-EK-foc-C}\\
\partial \mathcal{L}/\partial \ell_l &= U_l - \mu = 0, \label{eq:appB-EK-foc-l}\\
\partial \mathcal{L}/\partial \ell_f &= \lambda P_q F_L - \mu = 0, \label{eq:appB-EK-foc-f}\\
\partial \mathcal{L}/\partial \ell_o &= \lambda w_m - \mu = 0, \label{eq:appB-EK-foc-o}\\
\partial \mathcal{L}/\partial \ell_h &= \lambda (P_q F_L - w_m) - \mu m = 0. \label{eq:appB-EK-foc-h}
\end{align}

\paragraph{Shadow-wage divergence.} From (\ref{eq:appB-EK-foc-f}) and
(\ref{eq:appB-EK-foc-o}), the family-labor shadow wage equals the market wage
when an interior $\ell_o$ exists:
\begin{equation}
w_{f,i} \;\equiv\; P_q F_L \;=\; \mu / \lambda \;=\; w_m \quad \text{if } \ell_o > 0.
\label{eq:appB-EK-wf-interior}
\end{equation}
However, substituting into (\ref{eq:appB-EK-foc-h}),
\begin{equation}
\lambda(P_q F_L - w_m) \;=\; \mu m \;=\; \lambda w_m m
\quad \Longrightarrow \quad P_q F_L \;=\; w_m (1 + m).
\label{eq:appB-EK-MPL}
\end{equation}
The marginal product of farm labor must equal $w_m(1+m)$ rather than $w_m$
when hired labor is used; the supervision wedge $w_m \cdot m$ drives a gap
between the labor-market shadow wage on hired labor ($w_m + m \mu/\lambda$) and
the market wage. This violates AHM separability: the family's labor allocation
depends on whether and how much hired labor is used, which in turn depends on
the household's wealth and consumption choices.

\paragraph{Comparative statics --- Step 1.} The shadow value of time
$\mu/\lambda$ depends on the household's overall budget. Total differentiation
of (\ref{eq:appB-EK-foc-C})--(\ref{eq:appB-EK-foc-h}) with respect to $T_i$
yields, after lengthy but standard manipulations,
\begin{equation}
\frac{\partial (\mu/\lambda)}{\partial T_i} \;=\; \frac{\partial w_{f,i}}{\partial T_i} \;=\; f(U_{CC}, U_{Cl}, F_{LL}, m) \;\ne\; 0
\label{eq:appB-EK-step1}
\end{equation}
in general, with sign depending on the curvature of $U$ in $(C, \ell_l)$.

\paragraph{Comparative statics --- Step 2: $\partial \ell_o / \partial T_{SFFP}$.}
From (\ref{eq:appB-EK-foc-o}), $\mu/\lambda = w_m$ at interior $\ell_o$. The
subsidy $T$ relaxes the budget constraint, raising $\lambda$ (the marginal
utility of consumption falls less rapidly than the marginal utility of
leisure rises), which by the FOC system can cause $\ell_o$ to fall (income
effect on leisure: $\ell_o \downarrow$, $\ell_l \uparrow$) or rise (relaxing
the cash-constraint motive that drove supervision-cost hiring: $\ell_h \downarrow
\Rightarrow$ family time freed $\Rightarrow$ $\ell_o \uparrow$).

The off-farm-income response is therefore
\begin{equation}
\boxed{\frac{\partial \text{off\_farm\_income}_i}{\partial T_{SFFP}} \;=\; w_m \cdot \frac{\partial \ell_{o,i}}{\partial T_{SFFP}} \;\ne\; 0 \quad \text{under non-separability,}}
\label{eq:appB-EK-step2}
\end{equation}
recovering main-text equation~\ref{eq:EK-1}, with sign theoretically
indeterminate. Empirical estimation must resolve which of the two opposing
effects (income vs.\ supervision-relaxation) dominates.

\paragraph{Sufficient conditions.}
\begin{description}
\item[(SC6)] $\ell_h > 0$ for the supervision-wedge mechanism to bite. If the
  representative SFFP-eligible household uses no hired labor, the channel is
  inoperative. Empirically: FHES Wave~1 hired-labor share is approximately
  17--34\% across baseline-tenancy quintiles (low at high $s_0$, high at low
  $s_0$), supporting an interior $\ell_h > 0$ for the predominantly-tenant
  treated cohort.
\item[(SC7)] $U_{Cl}$ and $F_{LL}$ are bounded; the off-farm labor response is
  finite. Standard regularity.
\item[(SC8)] Interior $\ell_o > 0$ for the comparative-statics machinery (the
  alternative corner case is handled by checking the empirical share of
  households with $\ell_o = 0$; FHES Wave~1 shows fewer than 12\% of
  SFFP-eligible households at this corner).
\end{description}

\paragraph{Auxiliary status reminder.} The supervision channel is auxiliary to
\S B.1. Its rejection (falsification trigger F2 in the main text) does not
invalidate \S B.1's primary prediction. The contribution narrative under
$\hat\beta_4 = 0$ becomes: ``wealth-biased liquidity is the sole margin of
separability failure for Korean smallholders'' --- a tighter but still
publishable finding.

% -----------------------------------------------------------------------------
\subsection{Reduced-Form-to-Econometric-$\beta$ Mapping Table}
\label{sec:appB-mapping}
% -----------------------------------------------------------------------------

Table~\ref{tab:appB-mapping} expands the unified prediction table of the main
text (Table~\ref{tab:alpha3-predictions}) with a per-bandwidth grid: each
channel reports the empirical magnitude (P3b-2 estimates where available)
\emph{separately} across the three bandwidths T1 ($\pm 500$~m\textsuperscript{2}),
T2 ($\pm 1{,}000$~m\textsuperscript{2}), and T3 (MSE-optimal), the
operationalized falsification trigger per bandwidth, and the journal-grade
implication of falsification. Bandwidth specifications follow R-conventions \S 10
and are pre-locked alongside the four-bin $s_0$ partition (cf.\ \S B.1
post-(\ref{eq:appB-step3}) footnote).

\begin{landscape}
\begin{table}[ht]
\centering
\caption{Expanded reduced-form-to-econometric-$\beta$ mapping with per-bandwidth empirical fit and operationalized falsification triggers. Triggers tie to R-conventions \S 10 bandwidth grid and ADR-0002 outcome-hierarchy pre-lock.}
\label{tab:appB-mapping}
\small
\setlength{\tabcolsep}{4pt}
\begin{tabular}{p{2.0cm} p{1.6cm} p{0.9cm} p{1.0cm} p{3.6cm} p{3.4cm} p{4.8cm} p{2.0cm}}
\toprule
Channel & Outcome & Sign & Econ. $\hat\beta$ & Empirical fit by bandwidth (P3b-2) & Falsification trigger (per bandwidth) & Operational decision rule & Reframe \\
\midrule
Liquidity primary \#1 \citep{CarterOlinto2003_liquidity} & $A_{own}$ & $>0$ & $\hat\beta_1$ & T1: TBD; T2: $+1{,}089$ m\textsuperscript{2} pure tenant ($p\!=\!.033$); T3: $+408$ ($p\!=\!.003$) & $\hat\beta_1 \le 0$ at T2 \emph{and} T3 & Reject H$_1$ if $p(\hat\beta_1 > 0) > .10$ at both T2 and T3 & Null estimate \\
Liquidity gradient \citep{CarterOlinto2003_liquidity} & $A_{own} \times (1-s_{0,i})$ & $> 0$, mono in $1-s_{0,i}$ & $\hat\beta_2$ & T2 monotone across bins: $\{$ pure\_tenant $+1{,}089$ ($p\!=\!.033$); low\_owner $+410$ ($p\!=\!.051$); mixed $+393$ ($p\!=\!.063$); pure\_owner $0$ ref $\}$; T1: TBD; T3: TBD & F1 fires if $\hat\beta_2 \le 0$ at T2 \emph{and} four-bin point estimates fail weakly monotone in $(1-s_{0,i})$ at T2 (strict-AND boolean) & F1 fires if $p(\hat\beta_2 > 0) > .10$ at T2 \emph{and} four-bin point estimates not weakly monotone in $(1-s_{0,i})$ at T2 (Holm-corrected across the two conditions) & Universal income effect (drop wealth-bias) \\
Capital adj. (primary \#2) \citep{Kazukauskas2013_disinvestment} & op\_cost\_ex\_rent & $\le 0$ if $T_{SFFP}/\tau \!<\! 1$ & $\hat\beta_3$ & T1: $-4.02$M~KRW ($p\!=\!.055$); T2: $-3.22$M~KRW ($p\!=\!.079$); T3: TBD & $\hat\beta_3 > 0$ at \emph{any} bandwidth & Reconsider $\tau$ if $\hat\beta_3 > 0$ at T1 or T2 ($p < .10$); robustness range $\tau \in [20\text{M}, 40\text{M}]$~KRW & Reconsider $\tau$ calibration (\S\ref{sec:calibration}) \\
Labor (auxiliary) \citep{EswaranKotwal1986_supervision} & off\_farm\_income & $\ne 0$ under non-sep. & $\hat\beta_4$ & T1: TBD; T2: TBD; T3: TBD (queued for next cycle) & F2: $\hat\beta_4 = 0$ all bandwidths & F2 fires if $p(\hat\beta_4 \ne 0) > .10$ at all three bandwidths; informative, not rejecting (EK-1 sign indeterminacy per main-text \S\ref{sec:predictions}) & Drop supervision channel; primary survives \\
\midrule
\multicolumn{8}{l}{\textit{Ex-theory (B.1) --- aggregate-equilibrium implication, not a theory test:}} \\
Incidence (ex-theory) \citep{Floyd1965_incidence, AlstonJames2002_incidence} & unit\_rent\_price & $< 0$ via incidence closure & $\hat\beta_5$ & T1: $-48$ to $-130$ KRW/m\textsuperscript{2} across non-pure-owner bins; T2 pass-through $-11.1\%$; T3: TBD & N/A (not a theory test) & N/A & N/A \\
\bottomrule
\end{tabular}
\end{table}
\end{landscape}

\paragraph{Reading the table.} The two ``Liquidity'' rows correspond to the
primary contribution under ADR-0002 outcome-hierarchy ($A_{own}$ promoted to
primary \#1). Both must hold for the wealth-biased liquidity mechanism to be
supported; rejection of either row at the operational decision rule reframes
the paper as indicated in the rightmost column. The ``Capital adj.'' row is the
secondary primary outcome; positive $\hat\beta_3$ at T1 or T2 would force a
reconsideration of the calibrated adjustment threshold $\tau$ within the
robustness range $\tau \in [20\text{M}, 40\text{M}]$~KRW. The ``Labor'' row is
auxiliary: F2 is \emph{informative-but-not-rejecting} (EK-1 sign indeterminacy
per main-text \S\ref{sec:predictions}); $\hat\beta_4 = 0$ across all
bandwidths is consistent with supervision-cost margin being inoperative at the
SFFP scale but does not on its own falsify the channel. The ``Incidence'' row
is reported as evidence consistent with the model's empirical footprint but is
not a primary theoretical contribution (see main text \S 3.4 equilibrium rent
caveat).

\paragraph{Joint test --- F1 alone is load-bearing.} F1 alone is the
rejection-bearing falsification trigger for the AHM-extension framing. Under F1
rejection (i.e., $\hat\beta_2 \le 0$ at T2 or four-bin monotonicity broken at
T1 \emph{and} T2), the joint contribution collapses to a universal income
effect under separability and the paper repositions as a precise
developed-country null estimate, complementing
\citet{LaFaveThomas2016_markets}. F2 is supporting evidence on the
labor-imperfection margin and does not contribute to the joint test.

\bibliographystyle{plainnat}
\bibliography{../../Bibliography_base}

\end{document}
